\chapter{Provenance}
\label{app:provenance}

\section{Every result file this report cites}

Law 3 says that no number appears in this report unless it is reproducible from
a committed file under \code{experiments/results/} whose manifest names a commit
this repository can produce. Table~\ref{tab:provenance} is that chain, written
out.

\begin{center}
\small
\footnotesize
\begin{longtable}{p{.46\textwidth}p{.24\textwidth}ll}
\toprule
\textbf{Run directory} & \textbf{Engine} & \textbf{Commit} & \textbf{Clean} \\
\midrule
\endfirsthead
\toprule
\textbf{Run directory} & \textbf{Engine} & \textbf{Commit} & \textbf{Clean} \\
\midrule
\endhead
\bottomrule
\endfoot
\rpath{experiments/results/test/2026-08-26T21-15-35Z_p6-rules_6457ac7} & \texttt{rules} & \texttt{6457ac7} & yes \\  % results: experiments/results/test/2026-08-26T21-15-35Z_p6-rules_6457ac7
\rpath{experiments/results/test/2026-08-26T21-17-53Z_p6-ngram_6457ac7} & \texttt{ngram} & \texttt{6457ac7} & yes \\  % results: experiments/results/test/2026-08-26T21-17-53Z_p6-ngram_6457ac7
\rpath{experiments/results/test/2026-08-26T20-32-48Z_p6-llm_6457ac7} & \texttt{llm} & \texttt{6457ac7} & yes \\  % results: experiments/results/test/2026-08-26T20-32-48Z_p6-llm_6457ac7
\rpath{experiments/results/analysis/2026-08-26T21-20-11Z_p6-analysis_6457ac7} & \texttt{llm against rules against ngram} & \texttt{6457ac7} & yes \\  % results: experiments/results/analysis/2026-08-26T21-20-11Z_p6-analysis_6457ac7
\rpath{experiments/results/bench/2026-08-26T20-32-48Z_bench_6457ac7} & \texttt{llm against rules against ngram} & \texttt{6457ac7} & yes \\  % results: experiments/results/bench/2026-08-26T20-32-48Z_bench_6457ac7
\rpath{experiments/results/analysis/2026-08-26T02-53-59Z_analysis_336175f} & \texttt{rules against ngram} & \texttt{336175f} & yes \\  % results: experiments/results/analysis/2026-08-26T02-53-59Z_analysis_336175f
\rpath{experiments/results/analysis/2026-08-26T06-23-41Z_p5r-analysis_0d900ff} & \texttt{rules against ngram} & \texttt{0d900ff} & yes \\  % results: experiments/results/analysis/2026-08-26T06-23-41Z_p5r-analysis_0d900ff
\rpath{experiments/results/test/2026-08-26T02-50-54Z_rules_57d2aed} & \texttt{rules} & \texttt{57d2aed} & yes \\  % results: experiments/results/test/2026-08-26T02-50-54Z_rules_57d2aed
\rpath{experiments/results/test/2026-08-26T02-52-08Z_ngram_57d2aed} & \texttt{ngram} & \texttt{57d2aed} & yes \\  % results: experiments/results/test/2026-08-26T02-52-08Z_ngram_57d2aed
\rpath{experiments/results/test/2026-08-26T06-18-34Z_p5r-rules_0d900ff} & \texttt{rules} & \texttt{0d900ff} & yes \\  % results: experiments/results/test/2026-08-26T06-18-34Z_p5r-rules_0d900ff
\rpath{experiments/results/test/2026-08-26T06-21-03Z_p5r-ngram_0d900ff} & \texttt{ngram} & \texttt{0d900ff} & yes \\  % results: experiments/results/test/2026-08-26T06-21-03Z_p5r-ngram_0d900ff
\end{longtable}

\captionof{table}{Every run cited in this report, its engine, the commit its
manifest names, and whether the working tree was clean when it was produced. A
run marked dirty is not citable and the traceability job fails on a citation of
one.}
\label{tab:provenance}
\end{center}

The check is not a pattern match on the shape of a citation. It walks the whole
chain for every reference in every prose file and every LaTeX source: the path
must exist, a manifest must sit beside it or above it, that manifest must name a
commit this repository can produce an object for, and the manifest must not mark
the run dirty.

That distinction was learned the hard way. The job went red on a check that had
passed locally minutes earlier, because the clone had been made at the default
depth of one, so it held only the tip commit and every citation of a run taken in
an earlier commit failed with a message saying the commit does not resolve in
this repository. That sentence was true of the clone and false of the repository.
The failure was reproduced locally before anything was changed, because a
failure that is fixed without being reproduced is a failure that is guessed at.

\section{The result files that are here and are not cited}

Two earlier sets of test split runs are in this repository and nothing in this
report cites them.

The first pair was taken on a corpus whose test partition held one template per
family. Those results were correctly taken and are correctly reported as the
first, underpowered measurement. They are not edited, not moved and not deleted,
because result files are append only history and the honest description of them
is what they are rather than something to be tidied away.

The second pair was taken after the corpus repair and before the language model
existed. Chapter~\ref{ch:corpus} cites those two for exactly one purpose, which
is to show what the repair did to each engine, and
Table~\ref{tab:engines-moved} is that comparison with the caveat attached that
the two columns are not two measurements of the same thing.

\section{Version strings in the committed run logs}

A reader looking at the older run logs will find a tool version string that does
not match the release the results are associated with. That is consistent rather
than wrong: those runs were taken before the version bump, the engine
description records the version that was actually running, and the manifest's
commit resolves to a tree in which that string is correct.

Re-running them to make the string prettier is exactly the regeneration this
project forbids, so they stay, and this paragraph is the explanation.

\section{How to reproduce any number in this report}

\begin{enumerate}
  \item Clone the repository at the commit a manifest names.
  \item Regenerate the corpus from the seed recorded in the corpus manifest and
    the base year recorded beside it. The output is byte identical or the
    generator has a defect.
  \item Install from the committed lock file.
  \item Run the evaluation for the engine and split named in the manifest's
    command line, which is recorded verbatim.
  \item Compare the metrics document against the committed one.
\end{enumerate}

The language model engine is the one step that will not reproduce bit for bit,
and that is recorded rather than papered over: the serving runtime does not
guarantee determinism, temperature and seed are pinned and are recorded in the
run, and the engine description says in words that determinism is not claimed.
